\chapter{A working toolkit for Gaussian models}
\label{ch:gaussian-toolkit}
Gaussian conditioning silently appears in Kalman filters, Gaussian
processes, factor analysis, and several multivariate estimators.
It belongs before relying on their formulas.

\begin{definition}[Multivariate normal vector]
A vector $X\in\mathbb R^d$ is jointly Gaussian if every scalar
linear combination $a^\top X$ is normally distributed (possibly
degenerate). Its mean $\mu$ and covariance $\Sigma\succeq0$
characterize its distribution. When $\Sigma\succ0$, its density is
\[
 p(x)=(2\pi)^{-d/2}|\Sigma|^{-1/2}
 \exp\!\left[-\frac12(x-\mu)^\top\Sigma^{-1}(x-\mu)\right].
\]
\end{definition}
\begin{theorem}[Gaussian conditioning by a Schur complement]
Suppose $(X,Y)$ is jointly normal and partition its mean and
covariance into $\mu_X,\mu_Y$ and blocks $\Sigma_{XX}$,
$\Sigma_{XY}$, $\Sigma_{YY}$. If $\Sigma_{YY}$ is invertible,
\[
\begin{aligned}
\E[X\mid Y=y]&=\mu_X+\Sigma_{XY}\Sigma_{YY}^{-1}(y-\mu_Y),\\
\Cov(X\mid Y)&=\Sigma_{XX}
 -\Sigma_{XY}\Sigma_{YY}^{-1}\Sigma_{YX}.
\end{aligned}
\]
\end{theorem}
\begin{proof}
Define $R=X-\mu_X-\Sigma_{XY}\Sigma_{YY}^{-1}(Y-\mu_Y)$.
Then $\Cov(R,Y)=0$. Joint Gaussianity makes $R$ and $Y$
independent. Expanding $\Cov(R)$ gives the Schur complement.
\end{proof}

\begin{proposition}[Wishart sample-covariance model]
If $Z_1,\ldots,Z_T$ are independent $\mathcal N_d(0,\Sigma)$
with known zero mean, then $\sum_{t=1}^T Z_tZ_t^\top$
has a Wishart distribution with scale $\Sigma$ and $T$
degrees of freedom. The usual centered sample covariance instead
involves $T-1$ degrees of freedom for independent Gaussian samples.
\end{proposition}
This distributional assumption is the model behind classical
random-covariance results; the same formula does not describe
arbitrarily correlated neuronal time bins.

\begin{proposition}[Multivariate Stein identity]
For $X\sim\mathcal N_d(\mu,\Sigma)$ and a differentiable
function $f$ with appropriate integrability and boundary decay,
\[
\E[(X-\mu)f(X)]=\Sigma\,\E[\nabla f(X)].
\]
\end{proposition}
\begin{proof}
For a nondegenerate Gaussian density $p$,
$\nabla p(x)=-\Sigma^{-1}(x-\mu)p(x)$.
Integrate $\nabla(fp)$ and use the vanishing boundary term.
\end{proof}

\section{Gaussian KL divergence}
For full-rank Gaussian distributions
$P=\mathcal N(\mu_0,\Sigma_0)$ and
$Q=\mathcal N(\mu_1,\Sigma_1)$ in dimension $d$,
\[
D_{\mathrm{KL}}(P\|Q)=\frac12\left[
\operatorname{tr}(\Sigma_1^{-1}\Sigma_0)
+(\mu_1-\mu_0)^\top\Sigma_1^{-1}(\mu_1-\mu_0)
-d+\log\frac{\det\Sigma_1}{\det\Sigma_0}\right].
\]
A proof expands $\E_P[\log p(X)-\log q(X)]$, then
uses the covariance trace identity. The formula requires
positive-definite covariance matrices; singular distributions
need measure-theoretic treatment of their supports.

\begin{example}[Neural latent variables]
A Kalman measurement update is repeated conditioning of a joint
Gaussian state--observation vector. Its familiar gain matrix is
therefore a Schur-complement calculation, not a property of
all nonlinear or non-Gaussian observation processes.
\end{example}
% --- expanded edition (October 2026) ---
\section{Linear-Gaussian models}
Most Gaussian computations in this book are instances of one pattern.
\begin{proposition}[Linear-Gaussian marginal and posterior]
Let $z\sim\Normal(m,P)$ and $y\mid z\sim\Normal(Hz,R)$ with $P,R\succ0$.
Then $y\sim\Normal(Hm,HPH^\top+R)$ and
\[
z\mid y\sim\Normal\bigl(m+K(y-Hm),\ (I-KH)P\bigr),
\qquad K=PH^\top(HPH^\top+R)^{-1}.
\]
\end{proposition}
\begin{proof}
The vector $(z,y)$ is jointly Gaussian with $\Cov(z,y)=PH^\top$ and
$\Cov(y)=HPH^\top+R$. Apply the conditioning theorem of this chapter.
\end{proof}
The Kalman filter (Result 88) applies this proposition once per time
step; factor analysis applies it with $z$ the latent factors and $y$ the
population response; Gaussian-process regression applies it with $z$ the
function values at test points.

\paragraph{Information form.}
Writing $\Lambda=P^{-1}$ (the precision), the posterior precision is
$\Lambda+H^\top R^{-1}H$ and the posterior mean solves
$(\Lambda+H^\top R^{-1}H)\hat z=\Lambda m+H^\top R^{-1}y$. In this form
independent observations add precision, which explains why combining
neurons improves a population estimate and why perfectly redundant
neurons (identical rows of $H$ with correlated noise) add less than
their number suggests.

\section{Gaussian-process regression of a firing-rate function}
Let $f\sim\mathcal{GP}(0,k)$ with squared-exponential kernel
$k(t,t')=s^2\exp(-(t-t')^2/(2\ell^2))$, observed with noise as
$y_i=f(t_i)+\epsilon_i$, $\epsilon_i\sim\Normal(0,\sigma^2)$. For test
times $t_\ast$, the proposition gives
\[
\E[f_\ast\mid y]=K_{\ast}(K+\sigma^2I)^{-1}y,\qquad
\Cov(f_\ast\mid y)=K_{\ast\ast}-K_\ast(K+\sigma^2I)^{-1}K_\ast^\top.
\]
The length scale $\ell$ plays the role of a smoothing bandwidth, and
maximizing the marginal likelihood
$\log p(y)=-\tfrac12y^\top(K+\sigma^2I)^{-1}y-\tfrac12\log\det(K+\sigma^2I)-\tfrac n2\log2\pi$
chooses it from the data. The Gaussian observation model is only an
approximation for spike counts; GPFA and related methods apply it to
square-root or smoothed counts, while point-process GP methods require
non-conjugate approximations.

\section{Entropy and mutual information of Gaussians}
A $d$-dimensional Gaussian has differential entropy
$h=\tfrac12\log\bigl((2\pi e)^d\det\Sigma\bigr)$. For jointly Gaussian
$(X,Y)$,
\[
I(X;Y)=\tfrac12\log\frac{\det\Sigma_{XX}\det\Sigma_{YY}}{\det\Sigma},
\]
and in the scalar case $I=-\tfrac12\log(1-\rho^2)$. A correlation of
$0.1$ carries $0.005$\,nats; $0.9$ carries $0.83$\,nats. For a Gaussian input sent through an additive-noise channel with a
fixed noise covariance, Gaussian noise minimizes the mutual information.
Gaussian formulas computed from second-order statistics are therefore a
lower bound in that additive-noise setting, but not a general bound on
the information carried by nonlinear or non-Gaussian neural responses.

\begin{example}[Information of a population code under Gaussian noise]
For $y=s\,u+\epsilon$ with scalar stimulus $s\sim\Normal(0,\sigma_s^2)$,
tuning vector $u\in\R^N$, and $\epsilon\sim\Normal(0,\Sigma)$,
\[
I(s;y)=\tfrac12\log\bigl(1+\sigma_s^2\,u^\top\Sigma^{-1}u\bigr).
\]
The quantity $u^\top\Sigma^{-1}u$ is the linear Fisher information of the
population. Correlations along $u$ reduce it; correlations orthogonal to
$u$ do not. The formula follows from the matrix determinant lemma
applied to $\det(\Sigma+\sigma_s^2uu^\top)$.
\end{example}

\section*{Exercises for this chapter}
\begin{exercise}
Derive the information-form posterior from the covariance-form
posterior using the Woodbury identity.
\end{exercise}
\begin{exercise}
Show that two jointly Gaussian variables are independent iff they are
uncorrelated, and give a non-Gaussian counterexample.
\end{exercise}
\begin{exercise}
Compute the Gaussian KL divergence between $\Normal(0,1)$ and
$\Normal(\mu,\sigma^2)$ and examine its behaviour as $\sigma\to0$.
\end{exercise}
\begin{exercise}
Derive $I(s;y)$ in the population-code example.
\end{exercise}
